%---------------------------------------------------------------------------%
%->> Frontmatter
%---------------------------------------------------------------------------%
%-
%-> 生成封面
%-

\maketitle% 生成中文封面
\MAKETITLE% 生成英文封面
%-
%-> 作者声明
%-
\makedeclaration% 生成声明页
%-
%-> 中文摘要
%-
\intobmk\chapter*{摘\quad 要}% 显示在书签但不显示在目录
\setcounter{page}{1}% 开始页码
\pagenumbering{Roman}% 页码符号

财富水平如何影响经济决策，是神经经济学关注的核心问题。边际效用递减理论预测，随着财富积累，同等财富增量的主观价值下降，个体的选择偏好随之改变。这一理论在人类研究中已积累了大量实验证据；非人灵长类研究亦发现猕猴能够使用并积累代币等符号性财富，且其风险偏好随财富水平的变化而动态调整。然而，现有范式存在两方面局限。其一，固定的代币兑换阈值引入了外生参考点，在分析财富的神经表征及其对决策的影响时需额外建模。其二，不同选项始终在代币得失这一维度上比较，缺乏跨类别的价值权衡，无法体现财富作为通用媒介的属性。为此，本研究设计了赚取--消费范式，训练三只猕猴在赚取代币与花费代币换取即时奖励之间进行抉择。猕猴可在任意时刻自由消费，无需等待固定兑换阈值，从而弱化了外生参考点的干扰。同时，消费与赚取分属初级和次级强化物，构成跨类别的价值权衡，而非同一维度上的得失比较。实验结果表明，猕猴不仅能根据选项价值做出决策，还能依据当前财富水平动态调整偏好：随着财富积累，消费意愿增强，赚取意愿减弱。这一结果符合边际效用递减理论的预测。效用模型能够定量拟合上述偏好模式，且模型给出的选项价值差与决策时间高度相关------价值越接近的试次，猕猴的反应时间越长。在此过程中，前扣带皮层（anterior cingulate cortex, ACC）、眶额叶皮层（orbitofrontal cortex, OFC）与背外侧前额叶皮层（dorsolateral prefrontal cortex, DLPFC）均编码当前财富状态及财富依赖的选项价值。不同神经元偏好不同的财富水平，整个神经元群体覆盖了完整的财富范围；同时这种编码遵循韦伯-费希纳定律：神经元的财富水平调谐曲线在对数尺度而非线性尺度上对称，故低财富区间的表征精度高于高财富区间，与感知系统对数量差异的主观编码方式一致。在上述共性之外，三个脑区的功能亦存在差异：ACC对财富状态的编码最为突出，并在每次选择结果出现时动态更新财富表征，但不编码赚取选项的价值；OFC同时编码财富状态、财富依赖的选项价值以及基于物品的选择信号；DLPFC虽编码财富水平，却不参与财富状态的动态更新，主要负责基于动作的选择信号编码。这一分工模式表明，财富状态并非游离于决策系统之外的背景变量，而是深度嵌入前额叶的价值编码与决策形成过程。这些发现为理解财富塑造经济决策的神经机制提供了单神经元层面的直接证据。

\keywords{财富，赚取–消费，前扣带皮层，眶额叶皮层，背外侧前额叶皮层}% 中文关键词
%-
%-> 英文摘要
%-
\intobmk\chapter*{Abstract}% 显示在书签但不显示在目录

The influence of wealth on economic decision-making is a central question in neuroeconomics. The theory of diminishing marginal utility predicts that as wealth accumulates, the subjective value of equivalent wealth increments decreases, such that individuals' choice preferences change accordingly. This theory has received extensive experimental support in human studies; non-human primate research has further shown that macaques can use and accumulate symbolic wealth in the form of tokens, and that their risk preferences shift dynamically with wealth level. However, existing paradigms have two limitations. First, a fixed threshold at which accumulated tokens are automatically redeemed for reward introduces an exogenous reference point, which requires additional modeling in any analysis of the neural representation of wealth and its effects on decision-making. Second, options are always compared along a single dimension of token gain or loss, lacking cross-category value trade-offs and thus failing to capture wealth's role as a general medium of exchange. To address these limitations, we trained three macaques to perform a newly designed earning-spending paradigm in which they chose between earning tokens and spending tokens for immediate reward. Macaques were free to spend at any time without waiting for a fixed redemption threshold, thereby minimizing the influence of exogenous reference points. Spending and earning also involved primary and secondary reinforcers respectively, constituting a cross-category value trade-off rather than a gain-or-loss comparison within a single dimension. Macaques made decisions based not only on varying offer values but also on their current wealth states: as wealth accumulated, they became more willing to spend and less willing to earn, consistent with the predictions of diminishing marginal utility theory. A utility model quantitatively captured these preference patterns, and the model-derived value difference between options was strongly correlated with reaction time---trials with more similar values elicited longer responses. Three prefrontal subregions---the anterior cingulate cortex (ACC), orbitofrontal cortex (OFC), and dorsolateral prefrontal cortex (DLPFC)---all encoded current wealth state and wealth-dependent offer value. Individual neurons preferred distinct wealth levels, and the population collectively tiled the full wealth range; moreover, this encoding followed the Weber-Fechner law: neuronal tuning curves for wealth were symmetric on a logarithmic rather than linear scale, yielding higher representational precision at low wealth levels than at high ones, consistent with the logarithmic encoding of quantity differences in perceptual systems. Beyond these shared properties, the three regions exhibited functional differences. ACC most prominently encoded wealth state and updated its wealth representation at each choice outcome, but did not encode the offer value of earning options. OFC simultaneously encoded wealth state, wealth-dependent offer value, and good-based choice signals. DLPFC encoded wealth level but did not participate in the dynamic updating of wealth state, and primarily encoded action-based choice signals. This division of labor suggests that wealth state is not a background variable external to the decision system, but is deeply embedded in the prefrontal circuitry underlying value coding and decision formation. These findings provide single-neuron-level evidence for understanding the neural mechanisms by which wealth shapes economic decision-making.

\KEYWORDS{Wealth, Earning-spending, Anterior cingulate cortex, Orbitofrontal cortex, Dorsolateral prefrontal cortex}% 英文关键词

\pagestyle{enfrontmatterstyle}%
\cleardoublepage\pagestyle{frontmatterstyle}%

%---------------------------------------------------------------------------%
